% GENERATED FILE. Edit canonical lessons, then run npm run generate:books.
% canonical-source-hash: fnv1a64:eeb1476d2860d411
% canonical-lessons: ZH-C358-limao, ZH-C358-zhenzheng, ZH-W358-zhen-true, ZH-C358-zhenxin, ZH-C358-qianxu, ZH-C358-mianzi

\chapter{Real Manners}
\label{ch:zh-real-manners}
\hlchaptermodality{358}

\begin{chapteropening}
\textbf{By the end of this chapter:} \emph{I can talk about good manners, being sincere and modest, and saving face.}

Complete the last lesson of chapter 358: I can talk about good manners, being sincere and modest, and saving face.
\end{chapteropening}

\section[lǐmào]{lǐmào — good manners; politeness}

\label{lesson:ZH-C358-limao}



\begin{quote}
Before the new one: say the Mandarin for to criticize, then the Mandarin for to praise.
\end{quote}

\subsection*{You'll want to know: lǐmào}
\emph{lǐmào} — \textquotedblleft{}good manners; politeness\textquotedblright{}.

\begin{quote}
\emph{zhège háizi hěn yǒu lǐmào} — This child has very good manners.
\end{quote}

Its characters come later; for now it is a word for the ear.

\begin{cousinweb}
\emph{lǐmào} is built from \emph{lǐ} and \emph{mào}, both inherited from Middle Chinese, \emph{lejX} and \emph{maewH} (Guangyun, in Baxter's transcription), which glosses the first, quoting an older dictionary, 'that by which the spirits are served and blessing is brought', and gives the second only as an older spelling. How the pair came to mean manners is not settled.
\end{cousinweb}

\begin{grammarlens}[title={What you've built}]
One more word for giving reasons.
\end{grammarlens}

\subsection*{Your turn}
\begin{itemize}
\raggedright
  \item \emph{Say it:} \emph{lǐmào}
  \item \emph{Say it:} \emph{lǐmào}, once more
  \item \emph{Say it:} \emph{zhège háizi hěn yǒu lǐmào} — This child has very good manners.
  \item \emph{Recall:} say the Mandarin for to criticize, then the Mandarin for to praise, then say \emph{lǐmào} again
\end{itemize}

\begin{tcolorbox}[breakable,colback=teal!4,colframe=teal!35!black,title={Before you move on}]
What does lǐmào mean? (\textquotedblleft{}Good manners; politeness\textquotedblright{}.) Say it once more.
\end{tcolorbox}
\section[\texorpdfstring{\zh{真正} (zhēnzhèng)}{真正 (zhēnzhèng)}]{\zh{真正} (zhēnzhèng) — real; genuine}

\label{lesson:ZH-C358-zhenzheng}



\begin{quote}
Before the new one: say the Mandarin for to praise, then the Mandarin for good manners.
\end{quote}

\subsection*{You'll want to know: \zh{真正}}
\textbf{\zh{真正}} — \emph{zhēnzhèng} — \textquotedblleft{}real; genuine\textquotedblright{}.

\begin{quote}
\textbf{\zh{我觉得这是真正的中国菜}} \emph{wǒ juéde zhè shì zhēnzhèng de Zhōngguó cài} — I think this is real Chinese food.
\end{quote}

The new character, in pieces: \zh{真} — ten strokes: \zh{十} \emph{shí} on top, a narrow box with two bars inside it under that, a long bar across, and two short dots under it, left then right. The next lesson puts it in your hand.

\begin{cousinweb}
\textbf{\zh{真正}} (\emph{zhēnzhèng}) is built from \zh{真} 'true' and \zh{正} 'proper', both inherited from Middle Chinese, \emph{tsyin} and \emph{tsyengH} (Guangyun, in Baxter's transcription), which glosses the first 'true, as against false' and the second 'proper'.
\end{cousinweb}

\begin{grammarlens}[title={What you've built}]
One more word for giving reasons.
\end{grammarlens}

\subsection*{Your turn}
\begin{itemize}
\raggedright
  \item \emph{Say it:} \emph{zhēnzhèng}
  \item \emph{Say it:} \emph{zhēnzhèng}, once more
  \item \emph{Say it:} \emph{wǒ juéde zhè shì zhēnzhèng de Zhōngguó cài} — I think this is real Chinese food.
  \item \emph{Recall:} say the Mandarin for to praise, then the Mandarin for good manners, then say \emph{zhēnzhèng} again
\end{itemize}

\begin{tcolorbox}[breakable,colback=teal!4,colframe=teal!35!black,title={Before you move on}]
What does \zh{真正} mean? (\textquotedblleft{}Real; genuine\textquotedblright{}.) Say it once more.
\end{tcolorbox}
\section[\texorpdfstring{\zh{真} (zhēn)}{真 (zhēn)}]{\zh{真} — ten, a box, a bar and two dots}

\label{lesson:ZH-W358-zhen-true}



\begin{quote}
The word first, then the character inside it.

\begin{itemize}
\raggedright
  \item \emph{Recall:} say \emph{zhēnzhèng} — \textbf{R1}, one lesson back, before it decays
\end{itemize}

\begin{quote}
\zh{十} \zh{正}
\end{quote}

\emph{Read it:} each one aloud, with nothing said first — \emph{shí}, \emph{zhèng}
\end{quote}

\subsection*{Script — ten, a box, a bar and two dots}
\hlblockfigure{\detokenize{figures/ZH-W358-zhen-true-filmstrip.pdf}}{How \zh{真} is written, stroke by stroke}

\begin{quote}
\zh{真}
\end{quote}

\textbf{One piece, top to bottom.} Count it in four parts, not ten strokes.

\textbf{\zh{十} on top, two strokes}: a horizontal, then a vertical down through it.

\textbf{A narrow box, five strokes}: a vertical, a horizontal that turns down, two short horizontals inside, then one closing the foot of the box.

\textbf{A long horizontal} across the whole width under the box.

\textbf{Two short strokes at the foot}: one falling to the left, then one falling to the right.

2 + 5 + 1 + 2 = ten.

\subsection*{Your turn}
\begin{itemize}
\raggedright
  \item \emph{Write it:} \zh{真}, all ten, with the model in front of you
  \item \emph{Cover:} write it again from memory, and count the pieces, not the strokes
\end{itemize}

\begin{tcolorbox}[breakable,colback=teal!4,colframe=teal!35!black,title={Before you move on}]
What sits on top of \zh{真}? (\textbf{\zh{十}}.) How many bars inside the box? (\textbf{Two}, and one closing it.) What comes last? (\textbf{Two short strokes, left then right}.)

Source: \href{https://raw.githubusercontent.com/chanind/hanzi-writer-data/68d10a4b21150cae5e1ebbd223eed289cf32d90c/data/\%E7\%9C\%9F.json}{Hanzi Writer Data \zh{真}, pinned stroke paths and medians}.
\end{tcolorbox}
\section[\texorpdfstring{\zh{真心} (zhēnxīn)}{真心 (zhēnxīn)}]{\zh{真心} (zhēnxīn) — sincere; from the bottom of your heart}

\label{lesson:ZH-C358-zhenxin}



\begin{quote}
Before the new one: say the Mandarin for good manners, then the Mandarin for genuine.
\end{quote}

\subsection*{You'll want to know: \zh{真心}}
\textbf{\zh{真心}} — \emph{zhēnxīn} — \textquotedblleft{}sincere; from the bottom of your heart\textquotedblright{}.

\begin{quote}
\textbf{\zh{我真心觉得你说得有道理}} \emph{wǒ zhēnxīn juéde nǐ shuō de yǒu dàoli} — I honestly think you've got a point.
\end{quote}

\begin{cousinweb}
\textbf{\zh{真心}} (\emph{zhēnxīn}) is built from \zh{真} 'true' and \zh{心} 'the heart', both inherited from Middle Chinese, \emph{tsyin} and \emph{sim} (Guangyun, in Baxter's transcription), which glosses the first 'true, as against false'.
\end{cousinweb}

\begin{grammarlens}[title={What you've built}]
One more word for giving reasons.
\end{grammarlens}

\subsection*{Your turn}
\begin{itemize}
\raggedright
  \item \emph{Say it:} \emph{zhēnxīn}
  \item \emph{Say it:} \emph{zhēnxīn}, once more
  \item \emph{Say it:} \emph{wǒ zhēnxīn juéde nǐ shuō de yǒu dàoli} — I honestly think you've got a point.
  \item \emph{Recall:} say the Mandarin for good manners, then the Mandarin for genuine, then say \emph{zhēnxīn} again
\end{itemize}

\begin{tcolorbox}[breakable,colback=teal!4,colframe=teal!35!black,title={Before you move on}]
What does \zh{真心} mean? (\textquotedblleft{}Sincere; from the bottom of your heart\textquotedblright{}.) Say it once more.
\end{tcolorbox}
\section[qiānxū]{qiānxū — modest}

\label{lesson:ZH-C358-qianxu}



\begin{quote}
Before the new one: say the Mandarin for genuine, then the Mandarin for from the bottom of your heart.
\end{quote}

\subsection*{You'll want to know: qiānxū}
\emph{qiānxū} — \textquotedblleft{}modest\textquotedblright{}.

\begin{quote}
\emph{nǐ tài qiānxū le} — You're being too modest.
\end{quote}

Its characters come later; for now it is a word for the ear.

\begin{cousinweb}
\emph{qiānxū} is built from \emph{qiān} and \emph{xū}, both inherited from Middle Chinese, \emph{khem} and \emph{xjo} (Guangyun, in Baxter's transcription), which glosses the first 'respectful' and 'to yield', and the second 'empty'.
\end{cousinweb}

\begin{grammarlens}[title={What you've built}]
One more word for giving reasons.
\end{grammarlens}

\subsection*{Your turn}
\begin{itemize}
\raggedright
  \item \emph{Say it:} \emph{qiānxū}
  \item \emph{Say it:} \emph{qiānxū}, once more
  \item \emph{Say it:} \emph{nǐ tài qiānxū le} — You're being too modest.
  \item \emph{Recall:} say the Mandarin for genuine, then the Mandarin for from the bottom of your heart, then say \emph{qiānxū} again
\end{itemize}

\begin{tcolorbox}[breakable,colback=teal!4,colframe=teal!35!black,title={Before you move on}]
What does qiānxū mean? (\textquotedblleft{}Modest\textquotedblright{}.) Say it once more.
\end{tcolorbox}
\section[\texorpdfstring{\zh{面子} (miànzi)}{面子 (miànzi)}]{\zh{面子} (miànzi) — face; your standing in others' eyes}

\label{lesson:ZH-C358-mianzi}



\begin{quote}
Before the new one: say the Mandarin for from the bottom of your heart, then the Mandarin for modest about yourself.
\end{quote}

\subsection*{You'll want to know: \zh{面子}}
\textbf{\zh{面子}} — \emph{miànzi} — \textquotedblleft{}face; your standing in others' eyes\textquotedblright{}.

\begin{quote}
\textbf{\zh{他很要面子}} \emph{tā hěn yào miànzi} — He cares a lot about saving face.
\end{quote}

\begin{cousinweb}
\textbf{\zh{面子}} (\emph{miànzi}) is built from \zh{面} 'the face' and \zh{子}, a noun ending, both inherited from Middle Chinese, \emph{mjienH} and \emph{tsiX} (Guangyun, in Baxter's transcription), which glosses the first 'to face' and the second 'a child'. The second is said light here.
\end{cousinweb}

\begin{grammarlens}[title={What you've built}]
One more word for giving reasons.
\end{grammarlens}

\subsection*{Your turn}
\begin{itemize}
\raggedright
  \item \emph{Say it:} \emph{miànzi}
  \item \emph{Say it:} \emph{miànzi}, once more
  \item \emph{Say it:} \emph{tā hěn yào miànzi} — He cares a lot about saving face.
  \item \emph{Recall:} say the Mandarin for from the bottom of your heart, then the Mandarin for modest about yourself, then say \emph{miànzi} again
\end{itemize}

\begin{tcolorbox}[breakable,colback=teal!4,colframe=teal!35!black,title={Before you move on}]
What does \zh{面子} mean? (\textquotedblleft{}Face; your standing in others' eyes\textquotedblright{}.) Say it once more.
\end{tcolorbox}
